% ============================================================
%  4. ARQUITECTURA LÓGICA DE LA SOLUCIÓN  (Subdocumento 4.1)
%  Informe 1 — Arquitectura lógica
%  Detalle de módulos y capas, propia de la solución Puelche
% ============================================================
\chapter{Arquitectura Lógica de la Solución}
\label{sec:arqlogica}

Este capítulo presenta la arquitectura lógica de la solución, distinta de la explicación
del alcance y del diagrama de solución. Describe cada módulo y cada capa con el detalle
suficiente para comprender cómo se compone la plataforma, evidenciando el carácter
híbrido exigido en el Artículo 16° de las Bases Administrativas.

\section{Esquema general de la solución}

La solución se estructura en una \textbf{arquitectura de ocho capas} que organiza la
separación de responsabilidades desde la presentación hasta la observabilidad, con un
\emph{bus} de integración y eventos que desacopla los módulos de negocio. Cada capa es
propia de la solución planteada para Puelche, no un diagrama genérico.

\diagrama{arq_logica_capas.png}{Arquitectura lógica de ocho capas de la solución}{fig:arqlogica_capas}

\diagrama{arq_contexto.png}{Diagrama de contexto de la solución Puelche}{fig:arq_contexto}

\begin{table}[htbp]
\centering
\small
\caption{Arquitectura de ocho capas y su aplicación en la distribuidora}
\label{tab:capas}
\begin{tabularx}{\textwidth}{@{}L{4.2cm}X@{}}
\toprule
\textbf{Capa} & \textbf{Aplicación en distribuidora} \\
\midrule
Presentación &
Portal de administración (Talca/Concepción); app móvil de preventa; app de reparto;
app de recepción y \emph{picking} en bodega; autoatención de clientes, transportistas
y proveedores. \\
\addlinespace
Borde y exposición &
CDN + WAF gestionado para el portal público y las interfaces de programación expuestas
a proveedores y cadenas; terminación TLS 1.3. \\
\addlinespace
Puerta de enlace de servicios &
Punto único de publicación de las interfaces de pedidos, crédito, trazabilidad, entrega
e intercambio con cadenas. \\
\addlinespace
Servicios de negocio &
Módulos sin estado: preventa; planificación de rutas; gestión de bodega y \emph{picking};
trazabilidad y retiros; cadena de frío; entrega y prueba de entrega; cobranza y efectivo;
devoluciones y envases; costeo; reposición y compras. \\
\addlinespace
Integración y eventos &
Bus de mensajería que desacopla preventa, crédito, bodega, ruta, reparto y facturación;
colas con reintento para dispositivos sin señal. \\
\addlinespace
Datos &
Base transaccional (pedidos, stock, trazabilidad de lote, temperatura, prueba de entrega)
separada del almacén analítico (OTIF, \emph{fill rate}, costo de servir). \\
\addlinespace
Seguridad transversal &
Identidad federada, cifrado en tránsito y en reposo, cifrado a nivel de campo para datos
comerciales, de pago y de geolocalización. \\
\addlinespace
Observabilidad transversal &
Métricas, registros y trazas correlacionados de Talca, Concepción, las tres plataformas
de \emph{cross-docking} y los dispositivos de terreno. \\
\bottomrule
\end{tabularx}
\end{table}

\section{Modelo de dominio de negocio}

El modelo de dominio se organiza por agregados desacoplados por el \emph{bus} de eventos,
permitiendo que cada proceso evolucione de forma independiente. Los agregados principales
son: \textbf{Pedido}, \textbf{Stock}, \textbf{Lote}, \textbf{Ruta}, \textbf{Entrega},
\textbf{Temperatura/Excursión}, \textbf{Cobro}, \textbf{Envase retornable} y
\textbf{Costo de servir}. La separación transaccional/analítica garantiza que la captura
en el punto del hecho no degrade el desempeño operacional (RT-05.05/05.29).

\diagrama{arq_dominio_agregados.png}{Mapeo de módulos de negocio a agregados de dominio}{fig:arq_dominio_agregados}

\section{Decisiones lógicas clave}

\begin{itemize}
  \item \textbf{Estilo arquitectónico:} \emph{monolito modular} evolutivo sobre
        microservicios acotados para los módulos de mayor escalabilidad estacional
        (RT-02.02), balanceando mantenibilidad y elasticidad durante el \emph{peak} de
        septiembre.
  \item \textbf{Reserva de stock y conflictos:} reserva efectiva al sincronizar con el
        servidor central; en \emph{offline} el stock es indicativo y los conflictos se
        resuelven por \emph{timestamp} (Decisión D8), con deduplicación determinista
        (RT-02.06/02.07).
  \item \textbf{Operación desconectada:} las funciones críticas operan un turno completo
        (14 h) y el CD un mínimo de 24 h sin enlace (RT-03.10/03.12/03.13).
  \item \textbf{Integración con el ERP:} capa anticorrupción (RT-05.20) para convivir
        con el ERP y el WMS 2013 sin ``dos verdades'' (Restricción R-04).
  \item \textbf{Canal moderno:} mensajería electrónica estructurada sin integración
        punto a punto por cadena (RT-05.23).
  \item \textbf{Registros de decisiones de arquitectura} (ADR) documentados (RT-02.04).
\end{itemize}

\section{Identificación del punto único de falla principal}

El análisis identifica como componente de mayor saturación en el \emph{peak} de
septiembre el motor de sincronización del \emph{bus} de eventos entre los dispositivos de
terreno y la nube (RT-02.11). Este se dimensiona de forma elástica y se somete a pruebas
de carga estacionales para garantizar la indisponibilidad cero de la ventana de despacho
matutino (05:30--07:00). El detalle del dimensionamiento y la estrategia de resiliencia se
desarrolla en la arquitectura física (capítulo siguiente).

\vspace{4pt}
{\small\color{lafroxGrato}\pendiente{} Incorporar el esquema lógico (diagrama de
capas), el diagrama de contexto y el mapeo de módulos a los agregados de dominio en la
versión final.}
